\section{Uniform passage estimates}
\label{sec:direct-passage-estimates}

Two degeneracies govern the endpoint comparison. At exponentially small
incoming heights, the quadratic boundary coefficient can produce an
unbounded power-law multiplier in each endpoint derivative. Near a
boundary root, the drift and the leading curvature both tend to zero,
so an unweighted error estimate cannot preserve the curvature sign.
We retain the common tail multiplier explicitly, remove the
height-dependent term by a boundary-flow coordinate, and fix the phase
of both projected orbits using the same central height. After a fixed
positive boundary-time advance, the phase error carries arbitrarily
high powers of the drift. These estimates give
the curvature criterion of \cref{lem:curvature-criterion}, including
the exterior contribution, and the one-sided slope comparison in the
weak-root range of \cref{cor:direct-weak-root-comparison}.

In all passage estimates below, derivatives are taken in the stated state
variables, with parameters and all parameter-only scales and actions
fixed. Bounds are uniform on the indicated parameter regions, and each
finite order \(k\) is fixed before choosing the small parameter threshold.

\subsection{Weighted tail estimates}

We first estimate the map from an incoming height \(t\) to the fixed
section \(z=z_*\), keeping its power-law multiplier explicit.
Write \(\mathcal G_\sigma(U,z)=B(\sigma\sqrt{2U},z)
-q(\sigma\sqrt{2U})/z\), as in
\eqref{eq:physical-graph-equation-new}. Fix a compact interval of
positive incoming energies \(\mathcal U=U(t)/t\). For a sufficiently small
fixed \(z_*\), the conditions
\begin{equation}
\label{eq:direct-tail-smallness}
 \frac{|q_0|}{t}+\frac{|q_1|}{\sqrt t}+|q_2|
 \quad\hbox{sufficiently small},\qquad 0<t<z_*,
\end{equation}
put all those solutions in a slightly enlarged invariant cone
\(c_-z\le U(z)\le c_+z\). Indeed, on this cone,
\[
 \frac{|q(\sigma\sqrt{2U})|}{z}
 \le C\left(\frac{|q_0|}{z}+\frac{|q_1|}{\sqrt z}
                  +|q_2|+\sqrt z\right).
\]
Choose \(c_-<b_0/4\) and \(c_+>4\sup B\), containing the incoming
energies in their interior. Reduction of \(z_*\) and
\eqref{eq:direct-tail-smallness} makes
\(b_0/2<\mathcal G_\sigma<2\sup B\). The vector field points into the
cone along both boundary lines.
This proves continuation to \(z_*\), also for nearby incoming
energies needed for differentiation.

To make the coefficient of \(\sqrt{z_*}\) independent of later scaled
height thresholds, fix a common cone using coefficient bounds on a
preliminary box in \((\eta,z)\). Enlarge it once to contain the
moving-entrance energy \(1/(2m_{\rm land})\) used in the adaptive and
annular landings, with \(m_{\rm land}\) fixed before the coefficient
and height thresholds. A scaled entrance with a larger energy-to-height
ratio is first continued to this cone. Constants for this finite
continuation may depend on the thresholds; the coefficient
\(C_{\rm cone}\) remains independent of them.

For a fixed positive scaled entrance outside the common cone, put
\(z=\ell^2Z,\ U=\ell^2\mathcal U\),
with fixed \(Z_0>0\) and \(\mathcal U(Z_0)=\mathcal U_0\) in a compact positive
interval, and set
\[
 \epsilon_{\rm in}=\ell+|q_0|/\ell^2+|q_1|/\ell+|q_2|.
\]
Write \(B_0=B(0,0,p)\). On any fixed positive-height interval the graph
equation is
\(\mathcal U_Z=B_0+O_{C^k}(\epsilon_{\rm in})\) on a fixed enlarged
positive energy interval. Integration and the finite variational
equations give
\begin{equation}
\label{eq:direct-finite-common-cone-entry}
 \mathcal U(Z_1;\mathcal U_0)=\mathcal U_0+B_0(Z_1-Z_0)
                +O_{C^k}(\epsilon_{\rm in}).
\end{equation}
Choose \(Z_1\) large after the entrance interval, and then
\(\epsilon_{\rm in}\) small; the ratio \(\mathcal U(Z_1)/Z_1\) lies
in the common cone. The constants in this finite continuation may
depend on \(Z_0,Z_1\) and the entrance interval. Its derivative with
respect to \(\mathcal U_0\) is \(1+O(\epsilon_{\rm in})\), and replacing the tail multiplier
based at \(\ell^2Z_0\) by that based at \(\ell^2Z_1\) introduces
only \((Z_1/Z_0)^{2q_2}=1+O(|q_2|)\), with \(Z_0\) and \(Z_1\) fixed.

\begin{lemma}[Weighted tail]
\label{lem:direct-weighted-tail}
For every fixed \(k\), set
\[
 \varpi_t(z)=(z/t)^{-2q_2},\qquad \varpi_t=\varpi_t(z_*),\qquad
 e(t)=\sqrt t+|q_1|/\sqrt t .
\]
If \(|q_2|<1/8\) and \eqref{eq:direct-tail-smallness} holds, then
the centered endpoint map satisfies
\begin{equation}
\label{eq:direct-centered-tail}
 \frac{U(z_*;t\mathcal U)-U(z_*;t\mathcal U_\circ)}{t\varpi_t}
 =a_{t,p}(\mathcal U-\mathcal U_\circ)+R_{t,p}(\mathcal U),\qquad
 c<a_{t,p}<C,\quad \|R_{t,p}\|_{C^k}\le C_ke(t).
\end{equation}
Its inverse in these centered coordinates has uniformly bounded
derivatives through order \(k\). If the incoming energies lie in the
common cone, then
\begin{equation}
\label{eq:direct-tail-first-slope}
 \left|\log a_{t,p}\right|
       \le C_{\rm cone}\left(\sqrt{z_*}+\frac{|q_1|}{\sqrt t}\right).
\end{equation}
Here \(C_{\rm cone}\) depends only on the common cone and the
uniform bounds for \(B\), \(R_q\) and their first derivatives
on the preliminary normal-rate box. The same constant applies
for all admissible \(t,z_*\) in that box, independently of the
later coefficient and height thresholds. For another fixed cone
the same bound holds with a constant depending on that cone.
For a moving entrance \(U(t)=ct\), with fixed
\(c>4\sup B\), let \(U_{\rm end}(t)\) be its energy at \(z_*\).
Then
\begin{equation}
\label{eq:direct-moving-tail}
 \begin{gathered}
 U_{\rm end}'(t)=\varpi_t a(t),\qquad
 c_1<a(t)<C_1,\\
 |(t\partial_t)^j a(t)|+|(t\partial_t)^j\log a(t)|
 \le C_j\left(e(t)+\frac{|q_0|}{t}\right),\quad 1\le j\le k .
 \end{gathered}
\end{equation}
Here the cone is enlarged to contain \(c\) before the smallness
constants are chosen.
\end{lemma}

\begin{proof}
We separate the quadratic term in the first energy derivative from the
coefficients that control the higher variations. On the cone,
\begin{align}
\label{eq:direct-tail-coefficients}
 \mathcal G_U&=-2q_2/z+R_1,\qquad
 |R_1|\le C(z^{-1/2}+|q_1|z^{-3/2}),\\
 |\mathcal G_{U^j}|&\le C_j(z^{1/2-j}+|q_1|z^{-j-1/2}),
 \qquad j\ge2.\nonumber
\end{align}
The constant drift disappears on differentiation, and the constant in
the first line depends only on the stated cone and coefficient bounds.
Hence the first variation is
\begin{equation}
\label{eq:direct-tail-variation}
 J(z;t)=\varpi_t(z)
     \exp\int_t^z R_1(U(\zeta),\zeta)\,d\zeta .
\end{equation}
Its quotient by \(\varpi_t(z)\) is positive and bounded; integrating the
bound for \(R_1\) in \eqref{eq:direct-tail-coefficients} bounds its
logarithm by the right side of \eqref{eq:direct-tail-first-slope}.

To estimate the higher variations, factor out the same multiplier. Put
\(V_j=t^{-1}\partial_{\mathcal U}^jU=\varpi_t(z)W_j\).
Here \(W_1(t)=1\) and \(W_j(t)=0\) for \(j\ge2\).
After cancellation of the quadratic linear term, their equations are
\[
 (W_j)_z=R_1W_j+
 \sum_{m=2}^j t^{m-1}\varpi_t(z)^{m-1}\mathcal G_{U^m}
 \sum_{\substack{i_1+\cdots+i_m=j\\i_\ell\ge1}}
       c_{i_1,\ldots,i_m}W_{i_1}\cdots W_{i_m}.
\]
Every variation in the inner sum has order less than \(j\), so it remains
to bound the coefficients of these products. The two required integrals obey
\begin{align}
\label{eq:direct-weighted-integrals}
 t^{m-1}\int_t^{z_*}z^{1/2-m}(z/t)^{2|q_2|(m-1)}\,dz
       &\le C_m\sqrt t,\\
 |q_1|t^{m-1}\int_t^{z_*}z^{-m-1/2}
                       (z/t)^{2|q_2|(m-1)}\,dz
       &\le C_m|q_1|/\sqrt t.\nonumber
\end{align}
After \(z=t\zeta\), the first integral has decay margin
\(m-3/2-2|q_2|(m-1)\ge1/4\); the second has a larger margin.
Variation of constants and induction give bounded \(W_1\) and
\(W_j=O(e(t))\) for \(j\ge2\). Choose \(a_{t,p}\) as the first
normalized derivative at \(\mathcal U_\circ\) and integrate the second
derivative bound. Constructing one extra derivative if necessary
proves \eqref{eq:direct-centered-tail}. Positivity of the first
derivative gives the inverse bounds.

For the moving entrance, differentiation includes the change of initial
height as well as that of initial energy. It gives
\[
 U_{\rm end}'(t)=J(z_*;t)\,f_c(t),\qquad
 f_c(t)=c-\mathcal G_\sigma(ct,t)>c/2.
\]
Thus \(a(t)=f_c(t)J(z_*;t)/\varpi_t\). To justify its derivatives,
write \(\mathcal E_t=t\partial_t\), holding \(z\) fixed inside the flow.
The differentiated equations and
\eqref{eq:direct-weighted-integrals} give
\begin{equation}
\label{eq:direct-moving-flow-jets}
 |\mathcal E_t^jU(z;t)|\le C_jt\varpi_t(z),\qquad j\ge1.
\end{equation}
For \(j=1\) this is \(\mathcal E_tU=tJf_c\). To compute higher moving
boundary data, set \(\mathcal B_j(t)=(\mathcal E_t^jU)(t;t)\). Then
\[
 \mathcal B_0=ct,\qquad
 \mathcal B_j=t\frac{d\mathcal B_{j-1}}{dt}
       -t(\partial_z\mathcal E_t^{j-1}U)(t;t).
\]
In the local variables \(z=t\zeta,\ U=tW\), near
\((\zeta,W)=(1,c)\), the equation has bounded finite jets in
\(\sqrt t,q_1/\sqrt t,q_0/t,q_2\). The positive-order Euler
derivatives of these quantities are bounded under
\eqref{eq:direct-tail-smallness}, so \(\mathcal B_j=O_j(t)\).
The same triangular induction then proves
\eqref{eq:direct-moving-flow-jets}.

These moving-flow bounds allow us to differentiate the logarithm in
\eqref{eq:direct-tail-variation}. The lower-boundary terms are
\(O(e(t))\). At order \(j\), the chain rule gives a finite sum
of terms with \(m\) flow variations, \(1\le m\le j\).
By \eqref{eq:direct-weighted-integrals} with \(m+1\) in place of
\(m\), each is bounded by \(C_je(t)\).
Repeated boundary differentiation has the same estimate.
Finally \(\mathcal E_t^jf_c=O(e(t)+|q_0|/t)\) for \(j\ge1\).
In particular its quadratic contribution \(2q_2c\) is constant
in \(t\). These estimates bound the Euler derivatives of \(a\); since
\(a\) has a positive lower bound, they also bound those of \(\log a\).
This proves \eqref{eq:direct-moving-tail}.
\end{proof}

For \(|q_2|<1/8\) and \(0<t\le z_*\), the tail amplitude satisfies
\begin{equation}
\label{eq:direct-tail-amplitude}
 t\varpi_t=z_*^{1/4}t^{3/4}(t/z_*)^{1/4+2q_2}\le z_*^{1/4}t^{3/4}.
\end{equation}

When the scaled entrance height also tends to zero, we first continue the
graph to a fixed scaled height where the weighted-tail estimate applies.
The following lemma supplies this continuation and the resulting endpoint
expansion. Its additional symbols and those of its proof are local.

\begin{lemma}[Continuation from a small scaled height]
\label{lem:direct-small-height-continuation}
Suppose \(0<\ell\le\delta\) and
\(|q_0|/(\delta\ell^2)+|q_1|/(\delta\ell)+|q_2|/\delta\le C\).
Put \(z=\ell^2Z\), \(U=\ell^2\mathcal U\) and
\(K_\sigma(x)=-q(\sigma\ell x)/(\delta\ell^2)\). Let a graph enter at
\(Z=\delta\zeta_0\) with \(\mathcal U=\mathcal U_0\in[0,\mathcal U_{\max}]\), where
\(0<c\le\zeta_0\le C\) and \(K_\sigma(\sqrt{2\mathcal U_0})\ge k_0>0\).
For every fixed \(Z_1>0\) and sufficiently small \(\delta\), it continues
with positive flux to
\[
 \mathcal U(Z_1)=\mathcal U_0+B_0Z_1+O(\delta(1+|\log\delta|)).
\]
A large fixed \(Z_1\), chosen after \(\mathcal U_{\max}\) and before \(\delta\),
puts \(\mathcal U(Z_1)/Z_1\) in the common cone.
If \(\mathcal U_0=x^2/2\), with \(x\) in a fixed compact positive interval and
\(\|\zeta_0(x)\|_{C^3}\le C\), set
\(t=\ell^2Z_1\), \(E=\ell^2\varpi_t\) and
\[
 b_\sigma=z_*-U(z_*;t,B_0t),\qquad
 A_\sigma=\varpi_t^{-1}
       \left.\partial_uU(z_*;t,u)\right|_{u=B_0t}.
\]
Then the endpoint satisfies
\begin{equation}
\label{eq:direct-small-height-endpoint}
 I_\sigma(x)=b_\sigma-\tfrac12EA_\sigma x^2+Ee_\sigma(x),\qquad
 c<A_\sigma<C,\quad
 \|e_\sigma\|_{C^3}\le C\delta(1+|\log\delta|).
\end{equation}
For two branches with the same \(\ell,Z_1\),
\(|\log(A_-/A_+)|\le2C_{\rm cone}\sqrt{z_*}+C\delta\).
The constants depend only on the displayed bounds, the compact radius
interval and \(Z_1\).
\end{lemma}

\begin{proof}
We first prove continuation, then retain the radius derivatives needed
for the endpoint expansion. The coefficient bounds and \(\ell\le\delta\)
bound all radius derivatives
of \(K_\sigma\) on compact intervals. Choose
\(\mathcal U_1>\mathcal U_{\max}+2Z_1\sup B\), then \(\zeta_1>\sup\zeta_0\) so large that
\(|K_\sigma|\le b_0\zeta_1/2\) for \(0\le \mathcal U\le \mathcal U_1\).
In \(Z=\delta\zeta\), the graph equation is
\[
 \mathcal U_\zeta=\delta\{B(\sigma\ell\sqrt{2\mathcal U},\delta\ell^2\zeta)
                       +K_\sigma(\sqrt{2\mathcal U})/\zeta\}.
\]
Up to \(\zeta_1\), the energy changes by \(O(\delta)\) and the radius by
\(O(\sqrt\delta)\), preserving \(K_\sigma\ge k_0/2\) and positive flux.
For \(Z\ge\delta\zeta_1\), \(|\delta K_\sigma/Z|\le b_0/2\) while
\(\mathcal U\le \mathcal U_1\), so \(b_0/2\le \mathcal U_Z\le2\sup B\), and this energy interval
persists to \(Z_1\). Integrating \(B-B_0=O(\ell)\) and
\(\delta K_\sigma/Z\) gives the displayed continuation estimate.
Since \(B_0\) is interior to the common cone, large \(Z_1\) gives entry.

For the radius-dependent entrance, the energies stay bounded away from
zero, so the preceding continuation can be differentiated in the entrance
radius. The flow in \(\zeta\) has uniformly bounded finite jets;
composition with \((\zeta_0(x),x^2/2)\) gives an \(O_{C^3}(\delta)\)
change up to \(\zeta_1\). Thereafter the coefficient error and its energy
derivatives are bounded by \(C_j(\delta+\delta/Z)\). Integration of the
finite variational equations gives the same continuation estimate in
\(C^3\). At \(t\), the smallness expression in
\eqref{eq:direct-tail-smallness} and \(e(t)\) are \(O(\delta)\).
Apply \eqref{eq:direct-centered-tail} with center
\(\mathcal U_\circ=B_0\) to the incoming energy \(\mathcal U(Z_1;x)/Z_1\).
The coefficient of this centered map is \(A_\sigma\), the normalized center derivative,
and \(t\varpi_t=Z_1E\); this proves
\eqref{eq:direct-small-height-endpoint}.
Finally \eqref{eq:direct-tail-first-slope}, with
\(|q_1|/\sqrt t=O(\delta)\), gives the slope-ratio bound.
\end{proof}

\subsection{A coordinate following the boundary flow}

The passage to a fixed scaled height combines boundary motion with a
height-dependent term. Dividing by the boundary drift would be singular
at its roots. We instead construct a coordinate that absorbs the
height-dependent term and follows the boundary flow, uniformly through
those roots. On a fixed positive radius interval, write
the scaled height equation as
\begin{equation}
\label{eq:direct-height-system}
 \frac{dX}{d\log Y}=\delta F(X)+YG(X,Y),\qquad \delta>0,
\end{equation}
where \(F\) has bounded derivatives through the required order and
\(G(X,Y)=B_0/X+O_{C^k}(\delta)\), with \(B_0\) in a compact
positive interval.

\begin{lemma}[Boundary-flow coordinate]
\label{lem:direct-boundary-flow}
On a smaller positive radius interval and \(0\le Y\le Y_0\),
for every fixed finite \(k\) and all sufficiently small \(\delta\),
there is a uniformly \(C^k\) coordinate
\[
 \mathcal W(X,Y)=X+Y\mathcal W_1(X,Y),\qquad
 \mathcal W_X>0,\quad \mathcal W_Y<0,
\]
which satisfies \(\mathcal W_{\log Y}=\delta F(\mathcal W)\) along
\eqref{eq:direct-height-system}. Its inverse \(X=\chi(w,Y)\)
has
\begin{equation}
\label{eq:direct-projection-energy}
 \frac{\chi(w,Y)^2}{2}
   =\frac{w^2}{2}+B_0Y+\delta YE_\chi(w,Y),\qquad
 \|E_\chi\|_{C^k}\le C_k,\quad \chi_Y>0.
\end{equation}
The construction is uniform through simple or multiple zeros of \(F\).
The height bound \(Y_0\) depends only on the bounds for \(B_0\) and on the
lower ends of the two radius intervals; the threshold for \(\delta\) and
the constants \(C_k\) may depend on the whole interval.
\end{lemma}

\begin{proof}
We construct the coordinate by following the full passage backwards and
then undoing the boundary motion. Extend the coefficients smoothly to a
larger interval and then to bounded functions, preserving positive \(G\). Let \(\phi^u\)
be the flow of \(\delta F\). Follow
\eqref{eq:direct-height-system} backward for logarithmic time \(u\),
writing its radius as \(X_b(u)\) and its height as \(Y\e^{-u}\).
Define
\[
 \mathcal W(X,Y)=\lim_{u\to\infty}\phi^u(X_b(u)).
\]
The derivative of the expression inside this limit is
\begin{equation}
\label{eq:direct-projection-integrand}
 -Y\e^{-u}(\phi^u)'(X_b(u))G(X_b(u),Y\e^{-u}).
\end{equation}
To differentiate the limit, we compare growth of the flow derivatives
with the height decay in this integrand. For each fixed order,
boundary-flow derivatives grow at most as
a polynomial in \(u\) times \(\e^{C_k\delta u}\).
The backward equation is
\[
 \partial_uX_b=-\delta F(X_b)-Y\e^{-u}G(X_b,Y\e^{-u}).
\]
Its state variations have linear coefficient
\(O(\delta)+O(Y\e^{-u})\). In each higher variational equation,
the remaining terms contain lower variations; coefficients from
\(F\) carry \(\delta\), and those from \(G\) carry \(\e^{-u}\).
Variation of constants therefore gives the same
polynomial-times-exponential bounds at every fixed order.
Consequently
\[
 \mathcal W=X-Y\int_0^\infty \e^{-u}
       (\phi^u)'(X_b(u))G(X_b(u),Y\e^{-u})\,du.
\]
Every prescribed derivative of this integral is bounded by an
integrable function
\(C_k(1+u)^{N_k}\e^{-(1-C_k\delta)u}\), after reducing
\(\delta\). This proves convergence and \(\mathcal W=X+Y\mathcal W_1\)
with uniform derivative bounds for \(\mathcal W_1\).
The flow identity gives \(\mathcal W_{\log Y}=\delta F(\mathcal W)\).

The signs and the inverse energy expansion follow by comparison with
\(\delta=0\). At that parameter, the energy identity is
\(\mathcal W_0^2=X^2-2B_0Y\). Choose \(Y_0\) so that this backward
graph remains inside the positive interval; since \(\mathcal W_0\le X\), this
constrains only the lower ends of the intervals. Write
\(G=B_0/X+\delta g\), and interpolate by
\(\lambda\delta F,\ B_0/X+\lambda\delta g\), \(0\le\lambda\le1\),
using compatible extensions. A \(\lambda\)-variation has
forcing of size \(\delta\) and introduces only additional
polynomials in \(u\). Differentiating the preceding integral
therefore gives
\[
 \mathcal W-\mathcal W_0=\delta YR(X,Y),\qquad \|R\|_{C^k}\le C_k.
\]
The construction gives two additional derivatives whenever needed.
Since \(\mathcal W_{0,X}>c>0\) and
\(\mathcal W_{0,Y}<-c<0\), the same signs and uniform inverse bounds
hold for \(\mathcal W\). If \(\chi_0=(w^2+2B_0Y)^{1/2}\), the identity
\[
 \mathcal W_0(\chi,Y)-\mathcal W_0(\chi_0,Y)=-\delta YR(\chi,Y)
\]
and the positive lower bound on \(\mathcal W_{0,X}\) give
\(\chi-\chi_0=\delta YO_{C^k}(1)\).
Squaring proves \eqref{eq:direct-projection-energy}, while
\(\chi_Y=-\mathcal W_Y/\mathcal W_X>0\).
\end{proof}

\begin{remark}[Continuation of the projected state]
\label{rem:direct-projected-continuation}
If the projected orbit starts in a component where \(F>0\), it
cannot cross the first zero of \(F\) in finite boundary-flow time.
Its inverse is an outward graph of \eqref{eq:direct-height-system}, since
\begin{equation}
\label{eq:direct-inverse-flux}
 X_{\log Y}=\delta F(w)\chi_w+Y\chi_Y>0.
\end{equation}
This remains valid as the projected state approaches a double root.
\end{remark}

The projected state moves with bounded speed in the boundary-flow time
\(\delta\log Y\). This bounds the height at which a fixed radial advance
can start.

\begin{lemma}[Height before a radial advance]
\label{lem:direct-radial-advance}
In \cref{lem:direct-boundary-flow}, let \(|F|\le C_F\),
\(\delta\sup|E_\chi|\le B_0\) and \(X_1<X_2\). Suppose a solution of
\eqref{eq:direct-height-system} lies in the domain of \(\mathcal W\) at
the heights \(Y_1\le Y\le Y_2\), with radius at most \(X_1\) at \(Y_1\)
and at least \(X_2\) at \(Y_2\). If \(8B_0Y_2\le X_2^2-X_1^2\), then
\[
 Y_1\le Y_2\exp\Bigl(-\frac{X_2-X_1}{2C_F\delta}\Bigr).
\]
\end{lemma}

\begin{proof}
At a point of the solution with radius \(X\), height \(Y\) and projected
state \(w\), \eqref{eq:direct-projection-energy} gives
\(0\le X^2-w^2=2Y\{B_0+\delta E_\chi(w,Y)\}\le4B_0Y\).
Hence \(w\le X_1\) at \(Y_1\). At \(Y_2\) the hypothesis gives
\(w^2\ge(X_1^2+X_2^2)/2\), so \(w\ge(X_1+X_2)/2\). Since
\(w_{\log Y}=\delta F(w)\) along the solution, \(w\) changes by at most
\(C_F\delta\log(Y_2/Y_1)\) between the two heights, and this change is at
least \((X_2-X_1)/2\).
\end{proof}

\subsection{Contracting height coordinates and the exact phase}

The projected boundary-flow times must now be expressed in the common
central height \(h\). We first linearize the height equation over a
possibly multidimensional base flow. We then normalize the passage from
the central section and use that normalization to fix the phase of each
projected orbit.
The following height normalization uses contraction in backward
\(\tau\)-time.

\begin{lemma}[Contracting height coordinate]
\label{lem:direct-contracting-height-coordinate}
Consider a skew system
\[
 x_\tau=g(x),\qquad y_\tau=\mathfrak a(x)y+y^2\mathcal N(x,y)
\]
on a base domain \(\mathcal O\) and \(|y|\le2y_0\), with \(\mathfrak a,\mathcal N\) and
their \(C^{k+1}\) derivatives bounded for a prescribed finite
\(k\). Suppose \(\partial_y(\mathfrak a y+y^2\mathcal N)\ge2c_y>0\) there. Fix a set
of initial base points \(\mathcal O_0\subset\mathcal O\), and suppose
that for \(x\) in a neighborhood of \(\mathcal O_0\) the backward base
flow \(x_b(t;x)\) stays in \(\mathcal O\) with uniform bounds
\[
 |D_x^j x_b(t;x)|\le C_j(1+t)^{M_j}\e^{jc_x t},
 \qquad 1\le j\le k+1,
\]
where \((k+1)c_x<c_y\). After reducing \(y_0\), the
function
\[
 \Theta(x,y)=-\int_0^\infty
      y_b(t;x,y)\mathcal N(x_b(t;x),y_b(t;x,y))\,dt
\]
is uniformly \(C^{k+1}\) on \(\mathcal O_0\times[-y_0,y_0]\), is divisible
by \(y\) with a bounded
\(C^k\) quotient, and
\(\Psi(x,y)=y \e^{\Theta(x,y)}\) has positive height derivative and
uniformly bounded inverse height derivatives through order \(k\).
Along the system, \(\Psi_\tau=\mathfrak a(x)\Psi\).
\end{lemma}

\begin{proof}
The backward height satisfies \(|y_b(t)|\le |y|\e^{-2c_y t}\).
In its variational equations the highest height jet has a negative
coefficient bounded away from zero. Each pure base derivative of
\(\mathfrak a y+y^2\mathcal N\) contains the decaying height; the other forcings are
products of lower height jets and base-flow jets. Variation of
constants and induction therefore bound every derivative of the
integrand through order \(k+1\) by
\(C_j(1+t)^{m_j}\e^{-(2c_y-jc_x)t}\), after a
reduction of the decay margin. The stated inequality makes these
bounds integrable. Differentiation under the integral gives
\(\Theta=O(y)\) and its uniform derivatives; one additional derivative
gives the bounded quotient \(\Theta/y\). Thus \(\Psi_y>c>0\) on a
smaller height interval, with the asserted inverse bounds.
Finally the flow property gives \(\Theta_\tau=-y\mathcal N\), and hence
\(\Psi_\tau=\mathfrak a\Psi\).
\end{proof}

For the passage from the central section, we next normalize the height
at a fixed regular radius. On a regular central prefix, the scaled
height equation is
\begin{equation}
\label{eq:direct-source-prefix}
 y_x=\frac{xy}{\delta(K(x)+B(x,y)y)},\qquad
 y(0)=h/Q,\qquad K(x)\ge k_0>0
\end{equation}
on a fixed interval. Here \(Q>0\) depends only on parameters.
The required derivatives of \(K,B\) are uniformly bounded.
Use \(t=\delta^{-1}\mathcal A(x)\), where
\(\mathcal A(x)=\int_0^x u/K(u)\,du\). The height equation becomes
\[
 y_t=y+y^2\mathcal N_t(y),\qquad
 \mathcal N_t(y)=-B/(K+By).
\]
The inverse of \(x\mapsto t\) is continuous at zero; the normalization below
differentiates only the height.

\begin{lemma}[Exact central-height normalization]
\label{lem:direct-source-height}
At each fixed terminal radius \(x_i>0\), on a sufficiently
small fixed height interval there is a coordinate
\[
 \Psi_i(y)=y\exp \Theta_i(y),\quad \Theta_i(0)=0,\quad
 \Psi_i'(y)>c,\quad
 \|\Psi_i\|_{C^k}+\|\Psi_i^{-1}\|_{C^k}\le C_k
\]
such that
\begin{equation}
\label{eq:direct-prefix-action}
 \Psi_i(y(x_i))=\e^{\mathcal A(x_i)/\delta}h/Q .
\end{equation}
These bounds hold on arbitrarily long intervals of boundary time.
The same conclusion holds if
\(K(x)\asymp K_*+x\), \(K'(x)\ge c\), \(K_*/\delta\ge1\),
with \(0<K_*\le K_{\max}\) and fixed \(K_{\max}\),
and the other coefficient derivative bounds are uniform.
\end{lemma}

\begin{proof}
Let \(T_i=\mathcal A(x_i)/\delta\), and let \(x_b(u)\) and \(y_b(u;y)\)
be the radius and height at backward time \(u\).
Set
\[
 Z(u)=\e^uy_b(u;y),\qquad
 a_\delta(u)=\frac{\e^{-u}}{K(x_b(u))}.
\]
The backward height equation becomes
\[
 Z'=\frac{a_\delta(u)Z^2B(x_b(u),\e^{-u}Z)}
           {1+a_\delta(u)ZB(x_b(u),\e^{-u}Z)}
       =:\mathcal N_u(Z).
\]
The long-time estimate reduces to a uniform bound on this coefficient
and its integral. In both cases of the lemma,
\[
 \sup_{0\le u\le T_i}a_\delta(u)
       +\int_0^{T_i}a_\delta(u)\,du\le C.
\]
For \(K\ge k_0\), this follows from \(a_\delta\le \e^{-u}/k_0\).
For \(K\asymp K_*+x\), let \(u_\delta\) be the backward
time to \(x_i/2\). The fixed terminal segment gives
\(u_\delta\ge c/\delta\) and \(K(x_b)\ge c_i>0\) before
that time. Afterwards \(K\ge cK_*\ge c\delta\), so
\[
 \sup_{u\ge u_\delta}a_\delta(u)
       +\int_{u_\delta}^{T_i}a_\delta(u)\,du
       \le C\delta^{-1}\e^{-c/\delta}.
\]
These bounds are independent of the positive lower endpoint
of the allowed \(K_*\)-range.

On a sufficiently small fixed \(Z\)-interval the denominator
in \(\mathcal N_u\) is at least \(1/2\), and
\[
 |\mathcal N_u(Z)|\le Ca_\delta|Z|^2,\qquad
 |\partial_Z\mathcal N_u(Z)|\le Ca_\delta|Z|,\qquad
 |\partial_Z^j\mathcal N_u(Z)|\le C_ja_\delta\quad(j\ge2).
\]
Gr\"onwall's inequality and the finite variational equations, with
the uniform integral bound for \(a_\delta\), give \(|Z(u;y)|\le2|y|\)
and uniform height derivatives of every prescribed finite order on a
smaller fixed interval. The first variation is positive and bounded
above and below. Set \(\Psi_i(y)=Z(T_i;y)\); then
\(y_b(T_i;y)=h/Q\) gives \eqref{eq:direct-prefix-action}, and
\[
 \Theta_i(y)=\log\frac{\Psi_i(y)}y
 =\int_0^{T_i}
   \frac{a_\delta ZB(x_b,\e^{-u}Z)}
        {1+a_\delta ZB(x_b,\e^{-u}Z)}\,du .
\]
The same bounds, with one extra derivative, give \(\Theta_i=O(y)\)
with uniformly bounded derivatives. Thus \(\Psi_i'(0)=1\), and the
inverse function theorem gives the asserted inverse bounds.
\end{proof}

The preceding normalization fixes the terminal radius. We next control
the dependence of the central height on that radius.
Suppose \(x\) ranges in a fixed compact positive
interval with \(K\ge k_0\), and fix a small terminal height \(y_f>0\).
Then the corresponding central height satisfies
\begin{equation}
\label{eq:direct-terminal-radius-jets}
 \delta\log(h(x)/Q)=-\mathcal A(x)+\delta M_p(x),\qquad
 \|M_p\|_{C^k}\le C_k .
\end{equation}
Choose a smaller fixed \(x_0>0\) with a positive action gap to
all terminal radii. On \(x\ge x_0\), use time \(\tau=x/\delta\).
The system is
\[
 x_\tau=\delta,\qquad
 y_\tau=\mathfrak a(x)y+y^2\mathcal N(x,y),\qquad \mathfrak a(x)=x/K(x)>c.
\]
Extend the coefficients by fixed smooth cutoffs preserving the positive
rate, and reduce the height box to preserve its derivative margin.
The backward base \(x_b(u)=x-\delta u\) has bounded state derivatives,
so \cref{lem:direct-contracting-height-coordinate} applies with
\(c_x=0\). It gives a uniformly regular coordinate
\(\Phi(x,y)=y\e^{\Theta(x,y)}\), with the required state derivative
and inverse bounds, satisfying \(\Phi_\tau=\mathfrak a(x)\Phi\).
Compose it at \(x_0\) with
\cref{lem:direct-source-height}, whose coordinate there is
\(\Psi_0\). The unit
\[
 W_p(\xi)=\frac{\Phi(x_0,\Psi_0^{-1}(\xi))}{\xi},
 \qquad W_p(0)=1,
\]
has uniform derivative bounds, using one extra order at zero.
The composition gives
\[
 \Phi(x,y_f)=\e^{\mathcal A(x)/\delta}\frac hQ\,W_p(\xi_0),\qquad
 \xi_0=\e^{\mathcal A(x_0)/\delta}h/Q,\qquad W_p(0)=1.
\]
Rearrange this as
\[
 \xi_0W_p(\xi_0)=
 \e^{-[\mathcal A(x)-\mathcal A(x_0)]/\delta}\Phi(x,y_f).
\]
The positive action gap bounds the right side and its \(j\)-th
terminal-radius derivative by \(C_j\delta^{-j}\e^{-c/\delta}\).
The uniformly regular inverse of \(\xi\mapsto\xi W_p(\xi)\) preserves
these finite-order bounds. Taking logarithms gives
\eqref{eq:direct-terminal-radius-jets} with
\(M_p=\log\Phi(x,y_f)-\log W_p(\xi_0)\).

The central-height normalization and the boundary-flow coordinate now
give the phase of the projected orbit. At a fixed regular entry
\(x_L\), combine \cref{lem:direct-source-height,lem:direct-boundary-flow}.
In the notation \(Y_i=\delta y_i\), let
\(\xi=\e^{\mathcal A(x_L)/\delta}h/Q\), so that
\(Y_i=\delta\xi \Upsilon(\xi)\), where \(\Upsilon(0)=1\) and \(\Upsilon\) has
uniform derivative bounds. Because \(F(x_L)>c>0\), the
coordinate shift at entry corresponds to the boundary-flow time
\[
 \alpha_i(\xi)
 =\int_{x_L}^{\mathcal W(x_L,\delta\xi \Upsilon(\xi))}\frac{dx}{F(x)}
 =\delta O(\xi).
\]
The integral lies in a regular entry collar. At a fixed
terminal height \(Y_o\), the resulting formula is
\begin{equation}
\label{eq:direct-common-phase}
 w=\Phi_F^{\tau(T)}(x_L),\quad
 \tau=T+C_p+\delta\Xi(\xi),\quad
 T=-\delta\log(h/Q),\quad
 C_p=\delta\log(Y_o/\delta)-\mathcal A(x_L).
\end{equation}
Here \(\Xi=\alpha_i/\delta-\log \Upsilon\) vanishes at zero and has
uniform derivative bounds. Both phase origins are now fixed by the same
central height \(h\). Near a boundary root, \(F(w)\) may tend to zero;
the next lemma bounds the derivatives of the phase correction relative
to this drift.

\begin{lemma}[Phase error at a boundary root]
\label{lem:direct-weighted-phase}
Suppose the phase in \eqref{eq:direct-common-phase} satisfies
\(\tau\ge\tau_0>0\), its orbit has \(F>0\), and
\(F(x_L)\ge f_0>0\), \(|F'|\le L_F\).
For every fixed \(N,j\ge0\), writing \(\theta=\tau_T\), one has
\begin{equation}
\label{eq:direct-phase-weight}
 |\theta-1|+\sum_{\ell=1}^j|\partial_T^\ell\theta|
       \le\varepsilon_{N,j}(\delta)F(w)^N,\qquad
 \varepsilon_{N,j}(\delta)\longrightarrow0 .
\end{equation}
\end{lemma}

\begin{proof}
In terms of \(\xi\), \eqref{eq:direct-common-phase} gives
\[
 \xi=(Y_o/\delta)\exp[-\tau/\delta+\Xi(\xi)],\qquad
 \xi_T=-\xi/\delta,\qquad \theta=1-\xi\Xi'(\xi).
\]
The left side of \eqref{eq:direct-phase-weight} is at most
\(C_j\delta^{-j}\xi\). To compare this bound with the drift, use
\(d\log F/d\tau=F'\) along the boundary flow; hence
\(F(w)\ge f_0\e^{-L_F\tau}\).
After division by \(F(w)^N\), the error is bounded by
\[
 C_{N,j}\delta^{-j-1}
       \exp[-(1/\delta-NL_F)\tau]
 \le C_{N,j}\delta^{-j-1}\e^{-\tau_0/(2\delta)}.
\]
Reduction of \(\delta\) proves the assertion, uniformly up to
simple or multiple boundary roots.
\end{proof}

For a no-root bottleneck, we also need height derivatives that remain
uniform as its action grows. On a fixed interval
\([x_L,x_R]\) suppose
\(F>0\), with positive margins only in fixed collars of its
endpoints. Define
\[
 \alpha_j(Y)=\int_{x_j}^{\mathcal W(x_j,Y)}\frac{dx}{F(x)},\qquad
 \Psi_j(y)=y\exp[-\alpha_j(\delta y)/\delta],\quad j=L,R.
\]
These are uniformly regular height coordinates, since
\(\alpha_j(Y)=Y\mathcal A_j(Y)\). The exact identity is
\begin{equation}
\label{eq:direct-bottleneck-action}
 \Psi_R(y_R)=
 \exp\left(\frac1\delta\int_{x_L}^{x_R}\frac{dx}{F(x)}\right)
 \Psi_L(y_L).
\end{equation}
For a small outgoing height, the backward projected flow and
\eqref{eq:direct-inverse-flux} reconstruct the outward passage.
Let \(\mathcal A_L\) be the action of the preceding
regular prefix and write its inverse height coordinate as
\(y_L=P_p(\zeta)\), where
\(\zeta=\e^{\mathcal A_L/\delta}h/Q\). By
\cref{lem:direct-source-height}, \(P_p(0)=0\), \(P_p'(0)=1\),
and its derivatives through the prescribed order are uniformly
bounded on a fixed small interval. Put
\[
 \begin{gathered}
 \mathcal A_f=\int_{x_L}^{x_R}\frac{dx}{F(x)},\qquad
 \beta=\e^{-\mathcal A_f/\delta}\in(0,1],\\
 \xi=\e^{(\mathcal A_L+\mathcal A_f)/\delta}h/Q,
 \qquad \mathcal C_p=\Psi_L\circ P_p.
 \end{gathered}
\]
Thus \(\zeta=\beta\xi\) and
\(\mathcal C_p(x)=x\Upsilon_p(x)\), with \(\Upsilon_p(0)=1\).
Composition with the prefix gives
\begin{equation}
\label{eq:direct-bottleneck-inverse-composition}
 \begin{aligned}
 \Psi_R(y_R)&=\beta^{-1}\mathcal C_p(\beta\xi)
             =\xi \Upsilon_p(\beta\xi),\\
 y_R&=\Psi_R^{-1}\bigl(\xi \Upsilon_p(\beta\xi)\bigr).
 \end{aligned}
\end{equation}
The factor \(\beta^{-1}\) can be large, but differentiation of the
argument \(\beta\xi\) contributes powers of \(\beta\). With \(\beta\)
fixed, the derivatives in the expanded height \(\xi\) are
\[
 \partial_\xi^j[\beta^{-1}\mathcal C_p(\beta\xi)]
       =\beta^{j-1}\mathcal C_p^{(j)}(\beta\xi),\qquad j\ge1.
\]
Its first derivative has a uniform positive lower bound, and
all prescribed higher derivatives are bounded because \(\beta\le1\).
Choose a fixed \(\xi_* >0\) so that \(\beta\xi\) lies in the
prefix domain and \(\xi \Upsilon_p(\beta\xi)\) lies in the common
domain of \(\Psi_R^{-1}\) for \(0\le\xi\le\xi_*\).
The resulting map \(\xi\mapsto y_R\) has a uniformly positive
first derivative; its image contains a fixed small terminal-height interval,
and its inverse has uniformly bounded derivatives there.
Constructing one extra derivative before dividing by \(\xi\) gives
\(y_R=\xi V_p(\xi)\), where \(V_p(0)=1\), and both \(V_p\)
and its reciprocal have uniform derivatives through the prescribed
order. These bounds for derivatives in the expanded height \(\xi\)
are uniform in \(\mathcal A_f\).

\subsection{Endpoint derivatives and the logarithmic slope ratio}

We now transfer the drift estimates to the endpoints. Composing the
inverse boundary projection with the weighted tail gives the derivative
in the projected radius; the phase formula then converts it to a
derivative in \(T\). Keeping the drift factor through this conversion
allows comparison with the exterior contribution to the logarithmic
slope ratio.

\pagebreak[3]
\begin{lemma}[Endpoint after a boundary-flow passage]
\label{lem:boundary-flow-endpoint}
Let \(\ell>0\), and write \(|\eta|=\ell X\) and \(z=\ell^2Y\) on one
branch. Suppose that on a fixed positive radius interval the height
equation has the form \eqref{eq:direct-height-system} with small
parameter \(\delta\), and let \(w\) be the projected state of
\cref{lem:direct-boundary-flow}. Fix a small height \(Y_o>0\) and a
compact interval \(J_w\subset(0,\infty)\) on which
\eqref{eq:direct-projection-energy} holds at \(Y=Y_o\), and put
\[
 t=\ell^2Y_o,\qquad
 \varepsilon_o=\ell+\frac{|q_0|}t+\frac{|q_1|}{\sqrt t}+|q_2|,\qquad
 E=\ell^2\Bigl(\frac{z_*}{t}\Bigr)^{-2q_2}.
\]
For \(\delta\) and \(\varepsilon_o\) sufficiently small, the endpoint
\(I\) of the graph whose projected state at height \(t\) is
\(w\in J_w\) satisfies
\begin{equation}
\label{eq:boundary-flow-endpoint}
 -I_w=EA_p w\,U_p(w),\qquad c<A_p<C,\qquad
 U_p=1+O_{C^3}(\delta+\varepsilon_o),
\end{equation}
where \(A_p\) depends only on the parameters and \(U_p\) is a
dimensionless positive unit factor. The constants depend on
\(J_w\), \(Y_o\) and the fixed boundary-flow data. If the two branches
use the same \(\ell\) and \(Y_o\), then
\begin{equation}
\label{eq:boundary-flow-amplitude-ratio}
 \left|\log\frac{A_-}{A_+}\right|
 \le2C_{\rm cone}\sqrt{z_*}+O(\delta+\varepsilon_o),
\end{equation}
where \(A_\pm\) are the values of \(A_p\) on the negative and positive
branches, respectively.  The branch labels are fixed throughout the
passage; later contraction estimates use only the upper one-sided
consequence of this absolute comparison.
\end{lemma}

\begin{proof}
By \eqref{eq:direct-projection-energy}, the energy at height \(t\) is
\(U(t)=\ell^2\{w^2/2+B_0Y_o+\delta Y_oE_\chi(w,Y_o)\}\), and its
\(w\)-derivative is \(\ell^2w\{1+O_{C^3}(\delta)\}\) on \(J_w\). The
normalized energies \(U(t)/t\) lie in a fixed compact positive
interval. If this interval is not inside the common cone, continue the
graphs first by \eqref{eq:direct-finite-common-cone-entry} from \(Y_o\)
to a fixed larger \(Y_1\), whose energies lie in that cone. There the
error is \(O_{C^k}(\varepsilon_o)\) and the first derivative is
\(1+O(\varepsilon_o)\), since \(\varepsilon_{\rm in}\le C\varepsilon_o\)
for fixed \(Y_o,Y_1\). \Cref{lem:direct-weighted-tail} at
\(t_1=\ell^2Y_1\) gives a centered endpoint map with slope
\(t_1(z_*/t_1)^{-2q_2}a\), where \(c<a<C\) and, by
\eqref{eq:direct-tail-first-slope},
\(|\log a|\le C_{\rm cone}(\sqrt{z_*}+|q_1|/\sqrt{t_1})\), and with a
relative \(C^3\) error \(O(e(t_1))=O(\varepsilon_o)\). Since
\((t_1/t)^{2q_2}=(Y_1/Y_o)^{2q_2}=1+O(\varepsilon_o)\), the chain rule gives
\eqref{eq:boundary-flow-endpoint}. The factors \(\ell^2\),
\((z_*/t)^{-2q_2}\) and \((Y_1/Y_o)^{2q_2}\) are common to the two
branches, which gives the bound for \(A_-/A_+\).
\end{proof}

To compute the logarithmic rate in the common height coordinate,
consider an endpoint map on a projected-radius interval contained in a
fixed compact subset of \((0,\infty)\), of the form
\begin{equation}
\label{eq:direct-endpoint-unit}
 -I_w=EA_p w\,U_p(w),\qquad
 c<A_p<C,\quad U_p=1+O_{C^3}(\delta),\quad 0<E\longrightarrow0 ,
\end{equation}
as given by \cref{lem:boundary-flow-endpoint} when
\(\varepsilon_o=O(\delta)\); its amplitude satisfies
\(E\le C\ell^{3/2}\) by \eqref{eq:direct-tail-amplitude} with
\(t=\ell^2Y_o\). Here \(E=E(p)\) and \(A_p\) are independent of the
central height.
Use the phase \eqref{eq:direct-common-phase}, so that
\(w_T=\theta F(w)\), with \(\theta=\tau_T\) and \(F=K/w>0\).
Put \(m=hI_h=-\delta I_T\),
\(\mathfrak H=\partial_T\log m\).
In the identities below, primes on \(\theta,\mathfrak H\)
denote \(T\)-derivatives, and primes on \(K,F,\log U_p,\mathfrak H_0\) denote
\(w\)-derivatives:
\begin{align}
\label{eq:direct-rate-identities}
 m&=\delta EA_pK(w)\theta U_p(w),\\
 \mathfrak H&=\theta\mathfrak H_0(w)+\theta'/\theta,\qquad
 \mathfrak H_0=K'/w+F(\log U_p)',\nonumber\\
 \mathfrak H'&=\theta^2F\mathfrak H_0'+\theta'\mathfrak H_0
                 +\theta''/\theta-(\theta'/\theta)^2,\nonumber\\
 \mathfrak H_0'&=(wK''-K')/w^2+F'(\log U_p)'+F(\log U_p)''.\nonumber
\end{align}
By \cref{lem:direct-weighted-phase}, these identities imply
\begin{equation}
\label{eq:direct-rate-curvature}
 \mathfrak H'=F(w)\left\{\frac{wK''-K'}{w^2}+o(1)\right\},\qquad
 |I_T|+|I_{TT}|\le CEF(w).
\end{equation}
The errors are uniform up to a limiting boundary root, with the drift
factor retained in both estimates.

For the two endpoints use the same logarithmic height coordinate
\(T=-\delta\log(h/Q)\) in the logarithmic slope ratio
\eqref{eq:direct-completed-log-ratio}.
The exterior contribution is
\begin{equation}
\label{eq:direct-exterior-curvature}
 \partial_T^2[(\log\Pi_p')(I_-)]
       =(\log\Pi_p')''(I_-)I_{-,T}^2+(\log\Pi_p')'(I_-)I_{-,TT}.
\end{equation}
By \eqref{eq:direct-rate-curvature}, this term is
\(O(EF_-+E^2F_-^2)\).

The exterior contribution has the same drift weight as the negative
endpoint estimate and carries the vanishing amplitude \(E\). This leaves
the coefficients \(Q_\sigma=(wK_\sigma''-K_\sigma')/w^2\) of the two
drifts to determine the principal curvature of \(\mathcal L\).

\begin{lemma}[Curvature criterion]
\label{lem:curvature-criterion}
On an interval of central heights, suppose that both endpoints have the
form \eqref{eq:direct-endpoint-unit} in the common logarithmic height
coordinate \(T=-\delta\log(h/Q)\), that their phases satisfy the
hypotheses of \cref{lem:direct-weighted-phase}, and that the negative
projected state stays in a fixed compact subset of \((0,\infty)\). If
\(Q_-\le-c\) and \(Q_+\ge c\) there, then \(\mathcal L''<0\) for all
sufficiently small parameters, uniformly up to simple or multiple
boundary roots.
\end{lemma}

\begin{proof}
By \eqref{eq:direct-rate-curvature},
\(\mathfrak H_\sigma'=F_\sigma(w_\sigma)\{Q_\sigma+o(1)\}\). The exterior
term \eqref{eq:direct-exterior-curvature} is \(O(EF_-+E^2F_-^2)=o(1)F_-\),
since \(F_-\) is bounded and \(E\to0\). Hence
\[
 \mathcal L''=\mathfrak H_-'-\mathfrak H_+'+\partial_T^2[(\log\Pi_p')(I_-)]
 \le-\tfrac c2(F_-+F_+)<0.\qedhere
\]
\end{proof}

For a drift
\(K(w)=K(0)+K'(0)w+\tfrac12K''(0)w^2+w^3R_K(w)\), direct
differentiation gives
\begin{equation}
\label{eq:direct-curvature-coefficient-cubic}
 \frac{wK''-K'}{w^2}
 =-\frac{K'(0)}{w^2}+3R_K(w)+5wR_K'(w)+w^2R_K''(w).
\end{equation}

In the weak-root range, the first rate identity gives a one-sided slope
comparison. The next result expresses it in the common height coordinate,
retaining the change of scale from the positive branch's coordinate.

\begin{corollary}[Weak-root slope comparison]
\label{cor:direct-weak-root-comparison}
Let \(I_\pm(h)\) be the \(C^2\) endpoint maps on one interval of
central heights, with \(I_{\pm,h}>0\), and let \(a,\kappa>0\).
Suppose the positive endpoint map satisfies
\(-I_{+,w}=EA_pwU_p(w)\) with \(E>0\), \(c<A_p<C\) and
\(U_p=1+O_{C^2}(\kappa)\), where \(E,A_p\) are independent of the
central height. The radius
interval satisfies \(0<w_*\le w_a\le w<\alpha\le w^*<\infty\), with
fixed \(w_*,w^*\), and its exact phase uses \(\widehat T=-\kappa\log h\).
Assume the weighted-phase lemma
applies there, \(F=K/w>0\), the \(C^2\) bounds for \(K\) are
uniform, and \(K'(w)/w\le-c_0<0\).
In the common logarithmic height coordinate \(T=a\widehat T\), put
\(\mathfrak H_\sigma=\partial_T\log(hI_{\sigma,h})\).
Then, for sufficiently small \(\kappa\),
\begin{equation}
\label{eq:direct-weak-root-rate}
 \mathfrak H_+\le-\frac{c_0}{2a}.
\end{equation}
If also \(\mathfrak H_-\ge-C_n\), \(|I_{-,T}|\le\varepsilon_p\to0\),
and the \(C^2\) exterior map on these endpoint windows has positive slope with
\(|(\log\Pi_p')'|\le M\), then, for sufficiently small \(a\),
\(\mathcal L_T>0\) on this interval, and the return displacement has at
most two zeros there.
All bounds, the phase-lemma constants, and the convergence
\(\varepsilon_p\to0\) are uniform on a fixed preliminary parameter region
with small \(a\).
\end{corollary}

\begin{proof}
With \(\theta=\partial_{\widehat T}\tau\),
\eqref{eq:direct-rate-identities} gives
\[
 a\mathfrak H_+
 =\theta\{K'/w+F(\log U_p)'\}
                   +\theta_{\widehat T}/\theta.
\]
The bounded radius interval and the stated phase estimates make
the error \(o(1)\), proving \eqref{eq:direct-weak-root-rate}.
In the coordinate \(T\),
\[
 \mathcal L_T
 =(\log\Pi_p')'(I_-)I_{-,T}
                      +\mathfrak H_--\mathfrak H_+
 \ge \frac{c_0}{2a}-C_n-M\varepsilon_p>0.
\]
Fix the constants on the preliminary region, reduce \(a\),
and then choose the phase, tail, and parameter thresholds.
Since \(\mathcal L_T>0\), \eqref{eq:direct-completed-log-ratio} gives the
two-zero count.
\end{proof}
